\section{The finite vector model of a spherical triangle}
\label{sec:vector-model}

Let \(u,v,w\in\R^3\) be the unit vectors pointing to the vertices
\(A,B,C\), respectively.  The side \(a\) is opposite \(A\), hence it is the
shorter central angle between \(v\) and \(w\); cyclically for \(b,c\).  Put
\begin{equation}
 x=v\dotp w=\cos a,\qquad
 y=w\dotp u=\cos b,\qquad
 z=u\dotp v=\cos c.                                         \tag{3.1}
\end{equation}
Define the positive side sines
\begin{equation}
 s_a=\norm{v\crossp w},\qquad
 s_b=\norm{w\crossp u},\qquad
 s_c=\norm{u\crossp v}.                                    \tag{3.2}
\end{equation}
Since the sides lie strictly between \(0\) and \(\pi\), these norms are
positively apart from zero.  Lagrange's identity gives
\begin{equation}
 s_a^2=1-x^2,\qquad s_b^2=1-y^2,\qquad s_c^2=1-z^2.          \tag{3.3}
\end{equation}
Finally put
\begin{equation}
 \tau=\detm[u\ v\ w]=u\dotp(v\crossp w)>0.                 \tag{3.4}
\end{equation}
The sign of \(\tau\) is the orientation certificate.  Nondegeneracy supplies a
rational lower bound for it.

\subsection{The Gram identity}

\begin{lemma}[Gram determinant]
\label{lem:gram}
For unit \(u,v,w\),
\begin{equation}
 \tau^2=1+2xyz-x^2-y^2-z^2.                                 \tag{3.5}
\end{equation}
\end{lemma}

\begin{proof}
The square of a scalar triple product is the determinant of the Gram matrix:
\[
 \tau^2=
 \detm\!\begin{pmatrix}
 1&z&y\\ z&1&x\\ y&x&1
 \end{pmatrix}.
\]
Expanding this \(3\times3\) determinant gives (3.5).  In coordinates the claim
is a polynomial identity with integer coefficients, so its arithmetic proof
needs only distributivity and commutativity.
\end{proof}

A useful factorization is
\begin{equation}
 \tau^2=(1-y^2)(1-z^2)-(x-yz)^2
       =s_b^2s_c^2-(x-yz)^2,                               \tag{3.6}
\end{equation}
with two cyclic companions.

\subsection{Tangent directions and angle pairs}

At the vertex \(A\), the unit tangent direction toward \(B\) is
\begin{equation}
 t_{AB}=\frac{v-zu}{s_c},                                   \tag{3.7}
\end{equation}
and the unit tangent direction toward \(C\) is
\begin{equation}
 t_{AC}=\frac{w-yu}{s_b}.                                   \tag{3.8}
\end{equation}
Indeed, the numerators are orthogonal to \(u\), and their squared norms are
\(1-z^2=s_c^2\) and \(1-y^2=s_b^2\).  Define the trigonometric pair of the
interior angle \(A\) by
\begin{align}
 c_A&=t_{AB}\dotp t_{AC}
     =\frac{x-yz}{s_bs_c},                                  \tag{3.9}\\
 s_A&=u\dotp(t_{AB}\crossp t_{AC})
     =\frac{\tau}{s_bs_c}.                                  \tag{3.10}
\end{align}
Cyclically,
\begin{align}
 c_B&=\frac{y-zx}{s_cs_a},&s_B&=\frac{\tau}{s_cs_a},        \tag{3.11}\\
 c_C&=\frac{z-xy}{s_as_b},&s_C&=\frac{\tau}{s_as_b}.        \tag{3.12}
\end{align}

\begin{lemma}[Angle pairs are unit pairs]
\label{lem:angle-unit}
For each vertex,
\[
 c_A^2+s_A^2=1,
 \qquad c_B^2+s_B^2=1,
 \qquad c_C^2+s_C^2=1.
\]
Moreover \(s_A,s_B,s_C>0\).
\end{lemma}

\begin{proof}
Equation (3.6) says
\((x-yz)^2+\tau^2=s_b^2s_c^2\).  Divide by the positively separated
quantity \(s_b^2s_c^2\).  The other equations are cyclic.  Positivity follows
from \(\tau>0\) and the positive side sines.
\end{proof}

The constructive angle-chart procedure of Section~\ref{subsec:trig-names}
therefore produces unique angle names \(A,B,C\in(0,\pi)\) with
\(\cos A=c_A\), \(\sin A=s_A\), and cyclically.  Notice that the algebraic
proofs can work entirely with the pairs \((c_A,s_A)\); actual radian-valued
angle names are needed only for formulas involving sums, halves or area.

\subsection{A compact certificate for a spherical triangle}

For object-level formalization it is convenient to regard the following as the
primitive certificate:
\begin{equation}
 \mathcal T=(u,v,w;x,y,z;s_a,s_b,s_c;\tau;
             c_A,s_A,c_B,s_B,c_C,s_C),                     \tag{3.13}
\end{equation}
with the equations (3.1)--(3.12), unit-vector equations, and positive rational
lower bounds for \(s_a,s_b,s_c,\tau\).  Every equation is polynomial after
clearing the displayed positive denominators.  A checker for \(\mathcal T\)
uses only signed rational arithmetic, bounded square-root brackets and finite
comparisons.  It does not decide equality of arbitrary real names.

\begin{proposition}[Finite-coordinate sufficiency]
Every division-free spherical identity proved through
Section~\ref{sec:cotangent-special} follows from the polynomial equations in
\(\mathcal T\), the addition laws for a constructive trigonometric system, and
finitely many positive denominator witnesses.
\end{proposition}

\begin{proof}
The primary laws are immediate consequences of (3.9)--(3.12).  The remaining
identities are obtained by finite addition, multiplication, substitution and
cross multiplication.  Half formulas additionally use the positive branches
fixed by \(0<a,b,c,A,B,C<\pi\).  No compactness, choice, maximum principle or
uncountable construction enters.
\end{proof}
